% formulasheet.tex -- Module 8: Formula sheet explained.
% One \frow per formula card in ch1.tex ... ch5.tex, in card order, with the card's own
% notation and badge (default = on the sheet; [notonsheet]; [combine=...]).
% rearranged form of its [notonsheet] card, so it is listed here as "Learn it".
\unnumberedchapter[\faIcon{square-root-alt}]{Formula sheet explained}
\phantomsection\label{fs:chap}

This chapter lists every formula in Module~\modulenumber, chapter by chapter, in the same
order as the formula cards. Each row says when to use the formula. The last column tells
you about the official HSC formula sheet.
\textcolor{hscformula}{\bfseries\faCheck\ Yes}: it is printed on the sheet, so you are given
it in the exam. \textcolor{hscmistake}{\bfseries\faTimes\ Learn it}: it is not on the sheet,
so learn it or learn how to derive it (its formula card shows how).
\textcolor{hsctiptext}{\bfseries\faLink\ Combine}: it is not printed in this form, but you
can get it in one or two lines from items on the sheet.

\section*{\Cref{ch1:chap}: Origins of the Elements}
\begin{formulatable}
  \frow{$E = mc^2$}
    {Energy becomes matter (a pair needs $E_\mathrm{min} = 2mc^2$), or matter becomes
     energy (annihilation, fusion).}
  \frow[notonsheet]{$z = \dfrac{\Delta\lambda}{\lambda_0} = \dfrac{\lambda - \lambda_0}{\lambda_0}$\newline
      $z \approx \dfrac{v}{c}$\rule[-1.6ex]{0pt}{1ex}}
    {The redshift and speed of a galaxy moving much slower than light (not for sound).}
  \frow[notonsheet]{$v = H_0 d$ \qquad $t \approx \dfrac{1}{H_0}$}
    {Hubble's law for distant galaxies, and a rough age of the Universe ($H_0$ in
     \si{s^{-1}}).}
  \frow{$\lambda_\mathrm{max} = \dfrac{b}{T}$}
    {The temperature of a black body, such as a star or the CMB, from its peak wavelength.}
  \frow[notonsheet]{$\dfrac{L_A}{L_B} = 100^{\frac{M_B - M_A}{5}}$}
    {Compare two stars' luminosities from their magnitudes (5 smaller $=$ 100 times more).}
  \frow[combine={$E = mc^2$ and $\mathrm{u} = 931.5\ \mathrm{MeV}/c^2$}]%
    {$E = \Delta m\,c^2$\newline $E\,(\mathrm{MeV}) = \Delta m\,(\mathrm{u}) \times 931.5$}
    {The energy released by fusion in stars, from the mass lost.}
\end{formulatable}

\section*{\Cref{ch2:chap}: Structure of the Atom}
\begin{formulatable}
  \frow[notonsheet]{$v = \dfrac{E}{B}$}
    {Crossed $E$ and $B$ fields with the beam going straight (velocity selector).}
  \frow[notonsheet]{$\dfrac{q}{m} = \dfrac{E}{B^2 r}$}
    {Thomson: balance the fields, then bend the beam with the same $B$ alone.}
  \frow[notonsheet]{$\dfrac{q}{m} = \dfrac{2V}{B^2 r^2}$}
    {A school e/m tube: sped up from rest through $V$, then a circle at \ang{90} to $B$.}
  \frow[notonsheet]{$q = \dfrac{mgd}{V}$}
    {Millikan: an oil drop held still between horizontal plates.}
  \frow[notonsheet]{$q = ne$}
    {Check that a charge is a whole number of $e$, or count electrons gained or lost.}
  \frow{$F = \dfrac{1}{4\pi\varepsilon_0}\dfrac{q_1 q_2}{r^2}$}
    {The electric force between two small charges, such as an alpha particle and a nucleus.}
  \frow{$\sum m\vec{v}_\mathrm{before} = \sum m\vec{v}_\mathrm{after}$\newline
      $\sum\dfrac{1}{2}mv^2_\mathrm{before} = \sum\dfrac{1}{2}mv^2_\mathrm{after}$}
    {Collisions (Chadwick); the kinetic-energy equation is for elastic collisions only.}
\end{formulatable}

\section*{\Cref{ch3:chap}: The Quantum Atom}
\begin{formulatable}
  \frow{$E = hf$}
    {The energy of one photon, in joules, from its frequency (or the reverse).}
  \frow[combine={$E = hf$ and $v = f\lambda$}]{$E = \dfrac{hc}{\lambda}$}
    {Change between a photon's wavelength and its energy (not for particles with mass).}
  \frow{$\dfrac{1}{\lambda} = R\left(\dfrac{1}{n_\mathrm{f}^2} - \dfrac{1}{n_\mathrm{i}^2}\right)$}
    {The wavelength of any hydrogen line (Balmer: $n_\mathrm{f} = 2$; limit:
     $n_\mathrm{i} = \infty$).}
  \frow{$\lambda = \dfrac{h}{mv}$}
    {The wavelength of a moving particle with mass, below about $0.1c$.}
  \frow[notonsheet]{$2\pi r = n\lambda$}
    {Explain why only some Bohr orbits are allowed: a whole number of wavelengths fits.}
\end{formulatable}

\section*{\Cref{ch4:chap}: Properties of the Nucleus}
\begin{formulatable}
  \frow[notonsheet]{$\sum A_\mathrm{before} = \sum A_\mathrm{after}$\newline
      $\sum Z_\mathrm{before} = \sum Z_\mathrm{after}$}
    {Write or check any nuclear equation, and find a missing nucleus.}
  \frow{$N_\mathrm{t} = N_0 e^{-\lambda t}$}
    {The amount of one isotope left after any time $t$ ($\lambda$ and $t$ in one time
     unit).}
  \frow{$\lambda = \dfrac{\ln 2}{t_{\frac{1}{2}}}$}
    {Change between the half-life and the decay constant.}
  \frow[notonsheet]{$N_\mathrm{t} = N_0\left(\tfrac{1}{2}\right)^{n}$, \ $n = \dfrac{t}{t_{\frac{1}{2}}}$}
    {Quick answers when the time is a whole number of half-lives.}
  \frow[notonsheet]{$A = \lambda N$ \qquad $A_\mathrm{t} = A_0 e^{-\lambda t}$}
    {Activity in Bq from the number of nuclei ($\lambda$ in \si{s^{-1}}), and how it falls.}
  \frow{$E = mc^2$}
    {The energy of a nuclear reaction from the mass lost (in u: $\times\,931.5$ gives MeV).}
  \frow[notonsheet]{$\begin{aligned}[t]\Delta m &= (Zm_\mathrm{p} + Nm_\mathrm{n})\\
      &\quad - m_\mathrm{nucleus}\\ E_\mathrm{B} &= \Delta m\,c^2\end{aligned}$\rule[-1.6ex]{0pt}{1ex}}
    {How tightly a nucleus is held; divide by $A$ for the value per nucleon.}
  \frow[notonsheet]{$E = \sum E_\mathrm{B,after} - \sum E_\mathrm{B,before}$}
    {Energy from fission or fusion when binding energies are given instead of masses.}
\end{formulatable}

\section*{\Cref{ch5:chap}: Deep Inside the Atom}
\begin{formulatable}
  \frow{$\lambda = \dfrac{h}{mv}$}
    {The wavelength of a probe particle: more momentum shows smaller details.}
  \frow{$p_\mathrm{v} = \dfrac{m_0 v}{\sqrt{\left(1 - \dfrac{v^2}{c^2}\right)}}$}
    {The momentum of a particle faster than about $0.1c$ (use it in place of $mv$).}
  \frow[notonsheet]{$Q = q_1 + q_2 + q_3$ (baryon)\newline $Q = q_1 + q_2$ (meson)}
    {The charge of a hadron from its quarks (flip the sign for an antiquark).}
  \frow{$W = qV$}
    {The energy a charge gains in an accelerator gap ($n$ gaps from rest: $K = nqV$).}
  \frow[notonsheet]{$r = \dfrac{mv}{qB}$ \qquad $T = \dfrac{2\pi m}{qB}$}
    {A charge moving at \ang{90} to a magnetic field, well below $c$ (a cyclotron).}
\end{formulatable}

\section*{Other formulas used in this module}
These come from Modules 6 and 7, or appear in the text and on the formula cards.
\begin{formulatable}
  \frow{$v = f\lambda$}
    {The wave equation; for light, $c = f\lambda$ (\cref{ch1:chap,ch3:chap}).}
  \frow[combine={$E = hf$ and conservation of energy}]{$hf = E_\mathrm{i} - E_\mathrm{f}$}
    {The photon given out when an electron drops a level (\cref{ch3:chap}).}
  \frow[notonsheet]{$E_n = -\dfrac{\qty{13.6}{eV}}{n^2}$}
    {The energy of level $n$ in hydrogen (\cref{ch3:chap}).}
  \frow[notonsheet]{$mvr = \dfrac{nh}{2\pi}$}
    {Bohr's rule for the allowed orbits (\cref{ch3:chap}).}
  \frow{$d\sin\theta = m\lambda$}
    {A diffraction grating or a crystal (\cref{ch1:chap,ch3:chap}).}
  \frow{$P = \dfrac{\Delta E}{\Delta t}$}
    {Power: the energy a star or a reactor gives out each second. For a reactor, divide
     $P$ by the energy of one fission (in J) to get the fissions each second
     (\cref{ch1:chap,ch4:chap}).}
  \frow[combine={$P = \frac{\Delta E}{\Delta t}$ and $E = mc^2$}]{$\dfrac{\Delta m}{\Delta t} = \dfrac{P}{c^2}$}
    {The mass a star loses each second (\cref{ch1:chap}).}
  \frow{$E = \dfrac{V}{d}$ \qquad $\vec{F} = q\vec{E}$}
    {The field between parallel plates, and the force on a charge (\cref{ch2:chap}).}
  \frow{$F = qv_\perp B = qvB\sin\theta$}
    {The magnetic force on a moving charge (\cref{ch2:chap,ch5:chap}).}
  \frow{$F_\mathrm{c} = \dfrac{mv^2}{r}$}
    {The net force towards the centre; with $qvB$ it gives the radius
     (\cref{ch2:chap,ch5:chap}).}
  \frow{$K = \dfrac{1}{2}mv^2$}
    {Kinetic energy (\cref{ch2:chap,ch3:chap,ch5:chap}).}
  \frow[combine={$W = qV$ and $K = \frac{1}{2}mv^2$}]{$v = \sqrt{\dfrac{2qV}{m}}$}
    {The speed of a charge sped up from rest through $V$
     (\cref{ch2:chap,ch3:chap,ch5:chap}).}
  \frow[combine={$N_\mathrm{t} = N_0 e^{-\lambda t}$}]{$t = \dfrac{1}{\lambda}\ln\dfrac{N_0}{N_\mathrm{t}}$}
    {The time for a sample to fall from $N_0$ to $N_\mathrm{t}$ (\cref{ch4:chap}).}
  \frow[notonsheet]{$\lambda = \dfrac{h}{p_\mathrm{v}}$ \qquad $r = \dfrac{p_\mathrm{v}}{qB}$}
    {Wavelength and radius for a particle faster than about $0.1c$ (\cref{ch5:chap}).}
  \frow[notonsheet]{$f = \dfrac{1}{T} = \dfrac{qB}{2\pi m}$}
    {The fixed AC frequency of a cyclotron (\cref{ch5:chap}).}
\end{formulatable}

\section*{Data sheet values used in this module}
The exam gives you the data sheet. The values that Module~8 uses are below, exactly as
printed.

\begin{table}
  % unnumbered: a number here would print as ``Table 5.x'' (the chapter counter is 5)
  \caption*{Data sheet values used in Module~8.}
  \centering
  \small
  \begin{tabular}{@{}L{48mm}L{58mm}L{33mm}@{}}
    \toprule
    \colhead{Quantity (as on the sheet)} & \colhead{Value} & \colhead{Used in}\\
    \midrule
    Charge on electron & $q_\mathrm{e} = \qty{-1.602e-19}{C}$ & every chapter\\
    Mass of electron & $m_\mathrm{e} = \qty{9.109e-31}{kg}$ & Chapters 1, 2, 3, 5\\
    Mass of neutron & $m_\mathrm{n} = \qty{1.675e-27}{kg}$ & Chapters 2, 3\\
    Mass of proton & $m_\mathrm{p} = \qty{1.673e-27}{kg}$ & Chapters 1, 2, 3, 5\\
    Earth's gravitational acceleration & $g = \qty{9.8}{m.s^{-2}}$ & Chapter 2 (Millikan)\\
    Speed of light & $c = \qty{3.00e8}{m.s^{-1}}$ & Chapters 1, 3, 4, 5\\
    Electric permittivity constant & $\varepsilon_0 = \qty{8.854e-12}{A^2.s^4.kg^{-1}.m^{-3}}$
      & Chapter 2 (Coulomb)\\
    Planck constant & $h = \qty{6.626e-34}{J.s}$ & Chapters 1, 3, 5\\
    Rydberg constant & $R = \qty{1.097e7}{m^{-1}}$ & Chapter 3\\
    Atomic mass unit & $u = \qty{1.661e-27}{kg}$\newline $\phantom{u} = \qty{931.5}{MeV}/c^2$
      & Chapters 1, 2, 4\\
    1 eV & $\qty{1.602e-19}{J}$ & every chapter\\
    Wien's displacement constant & $b = \qty{2.898e-3}{m.K}$ & Chapter 1\\
    \bottomrule
  \end{tabular}
\end{table}

\begin{itemize}
  \item The sheet prints the electron's charge as negative. Use its size,
        $e = \qty{1.602e-19}{C}$, and find directions separately.
  \item \emph{Not} on the data sheet (the question gives them): $H_0$, the megaparsec,
        masses of nuclei in u, the \qty{13.6}{eV} of hydrogen, the density of oil, the
        laboratory wavelengths of spectral lines, the energy of one fission.
  \item With masses in u, multiply the mass lost by 931.5 to get MeV. Do not mix this with
        kg and $c^2$ in one answer: the two routes differ by about \qty{0.2}{\percent}.
\end{itemize}

\begin{mistakes}
\textbf{Same letter, different meaning.} Before you use a formula, check what each letter
stands for in it.\par\smallskip
\small\renewcommand\arraystretch{1.25}%
\begin{tabular}{@{}L{12mm}L{60mm}L{62mm}@{}}
  \colhead{Letter} & \colhead{Meaning 1} & \colhead{Meaning 2}\\
  $\lambda$ & wavelength: $E = \frac{hc}{\lambda}$, $\lambda = \frac{h}{mv}$, Rydberg
    & decay constant: $N_\mathrm{t} = N_0 e^{-\lambda t}$\\
  $E$ & energy: $E = hf$, $E = mc^2$ & electric field: $E = \frac{V}{d}$,
    $v = \frac{E}{B}$\\
  $e$ & elementary charge: $q = ne$, $e = \qty{1.602e-19}{C}$ & the number $e = 2.718$:
    $N_\mathrm{t} = N_0 e^{-\lambda t}$\\
  $A$ & mass number (protons $+$ neutrons) & activity: $A = \lambda N$, in \si{Bq}\\
  $N$ & number of neutrons, in $\Delta m = (Zm_\mathrm{p} + Nm_\mathrm{n}) - m_\mathrm{nucleus}$
    & number of nuclei: $N_\mathrm{t} = N_0 e^{-\lambda t}$\\
  $Z$, $z$ & $Z$: atomic number (protons) & $z$: redshift, $z = \frac{\Delta\lambda}{\lambda_0}$\\
  $T$ & temperature in kelvin: $\lambda_\mathrm{max} = \frac{b}{T}$ & period:
    $T = \frac{2\pi m}{qB}$\\
  $M$ & absolute magnitude of a star & $M_\odot$: the Sun's mass (and M is a spectral
    class)\\
  $m$ & mass, in \si{kg} & order of a bright line: $d\sin\theta = m\lambda$\\
  $n$ & level number: $E_n$, $2\pi r = n\lambda$, $mvr = \frac{nh}{2\pi}$ & a count:
    electrons in $q = ne$, half-lives in $\left(\frac{1}{2}\right)^{n}$, gaps in $K = nqV$\\
  $R$ & Rydberg constant, \qty{1.097e7}{m^{-1}} & a radius: $R_\odot$ (the Sun's radius),
    a cyclotron's dee\\
  u, $u$ & u: atomic mass unit, $\qty{1}{\amu} = \qty{931.5}{MeV}/c^2$ & u: the up quark;
    $u$: a velocity before a collision\\
  $V$, $v$ & $V$: voltage, in \si{V} & $v$: speed, in \si{m.s^{-1}}\\
\end{tabular}
\end{mistakes}

\begin{toolkit}{Change to SI units first}
Every formula in this module needs SI units, except the ``$\times\,931.5$'' rule (u to MeV).
\begin{itemize}
  \item nm to m: $\times 10^{-9}$. \ fm to m: $\times 10^{-15}$. \ cm to m: $\div 100$.
  \item eV to J: $\times\, \num{1.602e-19}$. \ MeV to J: $\times\, \num{1.602e-13}$.
        \ J to eV: $\div\, \num{1.602e-19}$.
  \item kV to V: $\times 1000$. \ MBq to Bq: $\times 10^6$.
  \item Hours, days or years to seconds: $\times 3600$, $\times \num{86400}$ or
        $\times\, \num{3.156e7}$. In $N_\mathrm{t} = N_0 e^{-\lambda t}$ any time unit
        works, if $\lambda$ and $t$ use the same one.
  \item \si{km.s^{-1}.Mpc^{-1}} to \si{s^{-1}}: divide by $\num{3.086e19}$ (the number of km
        in one Mpc).
\end{itemize}
\end{toolkit}
